\title{Разработка прототипа вопросно-ответного ассистента\\
  на основе Retrieval-Augmented Generation\\}

\taskdate{28.02.2026}

\projecttasks{%
  \projecttask{\bfseries\projecttasknum}{\bfseries Аналитическая часть}{}{}{}%
  \projecttask{\projectsubtasknum}%
  {Анализ проблемы поиска информации в корпоративной базе знаний.}%
  {Аналитический отчет}%
  {02.03.2026}{}%
  \projecttask{\projectsubtasknum}%
  {Обзор подходов к построению вопросно-ответных систем.}%
  {}%
  {06.03.2026}{}%
  \projecttask{\projectsubtasknum}%
  {Обоснование выбора Retrieval-Augmented Generation для первой MVP-версии.}%
  {Текст РСПЗ}%
  {10.03.2026}{}%
  \projecttask{\bfseries\projecttasknum}{\bfseries Теоретическая часть}{}{}{}%
  \projecttask{\projectsubtasknum}%
  {Рассмотрение основ больших языковых моделей.}%
  {}%
  {16.03.2026}{}%
  \projecttask{\projectsubtasknum}%
  {Рассмотрение векторного представления текста и векторного поиска.}%
  {}%
  {20.03.2026}{}%
  \projecttask{\projectsubtasknum}%
  {Описание базовой технологии Retrieval-Augmented Generation.}%
  {Алгоритм}%
  {24.03.2026}{}%
  \projecttask{\bfseries\projecttasknum}{\bfseries Инженерная часть}{}{}{}%
  \projecttask{\projectsubtasknum}%
  {Проектирование базового RAG-пайплайна.}%
  {}%
  {30.03.2026}{}%
  \projecttask{\projectsubtasknum}%
  {Проектирование архитектуры обработки документов.}%
  {}%
  {03.04.2026}{}%
  \projecttask{\projectsubtasknum}%
  {Проектирование архитектуры поиска и генерации ответа.}%
  {Диаграммы, схемы}%
  {10.04.2026}{}%
  \projecttask{\bfseries\projecttasknum}{\bfseries Технологическая и
    практическая часть}{}{}{}%
  \projecttask{\projectsubtasknum}%
  {Подготовка корпуса документов.}%
  {Отчеты JSON/JSONL, таблицы ПЗ}%
  {17.04.2026}{}%
  \projecttask{\projectsubtasknum}%
  {Реализация парсинга, чанкинга и индексации.}%
  {Исполняемые файлы}%
  {24.04.2026}{}%
  \projecttask{\projectsubtasknum}%
  {Реализация базового сценария вопрос-ответ и первичная проверка работоспособности.}%
  {Отчет проверки}%
  {30.04.2026}{}%
  \projecttask{\projectsubtasknum}%
  {Оформление ограничений MVP и направлений дальнейшего развития.}%
  {Текст ПЗ}%
  {05.05.2026}{}%
  \projecttask{\bfseries\projecttasknum}%
  {\bfseries Оформление пояснительной записки (ПЗ) и иллюстративного материала
    для доклада.}%
  {\bfseries Текст ПЗ, презентация}%
  {\bfseries 10.05.2026}{} }

\taskliterature{
\nocite{Lewis2020RAG}
\nocite{Vaswani2017Attention}
\nocite{Karpukhin2020DPR}
\nocite{Reimers2019SentenceBERT}
\nocite{QdrantDocs}
\nocite{SentenceTransformersDocs}
}

% # Подписи

% Для простановки подписи используются слдеющие команды:
% - простая подпись: \sign[<сдвиг>]{<масштаб>}{<FIO>}
% - подпись с датой: \signat[<сдвиг>]{<масштаб>}{<FIO>}{<дата>}
% где
% - <сдвиг> — необязательный сдвиг подписи по вертикали для правильного
%   расположения относительно строки
% - <масштаб> — число, используемое для масштарибования изображения подписи
% - <FIO> — имя файла с подписью, файл должен быть помещен в
%   img/signatures/FIO.png и иметь прозрачный фон
% - <дата> — дата, которая будет выведена под подписью

% ## Утверждение задания руководителем и студентом
\authortaskapproval{}%
\supervisortaskapproval{}%

% ## Утверждение РСПЗ руководителем, студентом и консультантом
\authorrspzapproval{}%
\supervisorrspzapproval{}%
\consultantrspzapproval{}%

% ## Оценка руководителя за РСПЗ
\supervisorrspzgrade{}%

% ## Утверждение ПЗ руководителем, студентом и консультантом
\authorpzapproval{}%
\supervisorpzapproval{}%
\consultantpzapproval{}%

% ## Оценка руководителя за ПЗ
\supervisorpzgrade{}%

%%% Local Variables:
%%% TeX-engine: xetex
%%% eval: (setq-local TeX-master (concat "../" (seq-find (-cut string-match ".*-1-task\.tex$" <>) (directory-files ".."))))
%%% End:
